\begin{statementbox}
\begin{center}\begin{tikzpicture}[font=\small,>=stealth,line join=round]
\draw[gray!40,rounded corners=3pt] (0,0) rectangle (14.2,8.3);\node at (7.1,7.700000000000001) {同过(1,1)，底数跨过1时，上下顺序互换};
\begin{scope}[shift={(1.6625,2.068181818181818)},x=1.8750000000000002cm,y=0.8363636363636363cm]\draw[->,gray] (-0.3,0)--(2.9,0) node[right]{$x$};\draw[->,gray] (0,-0.5)--(0,5) node[above]{$y$};\draw[blue,thick,solid,domain=0.03:2.15,samples=180] plot(\x,{(\x)^2});\draw[teal,thick,dashed,domain=0.03:2.6,samples=180] plot(\x,{\x});\draw[violet,thick,densely dotted,domain=0.03:2.6,samples=180] plot(\x,{sqrt(\x)});\fill[orange!85!black] (1,1) circle(2.2pt);\draw[gray,dashed] (1,0)--(1,4.6);\node[below=5pt] at (1,0){$1$};\end{scope}\node[anchor=west,align=left,text width=5.0cm] at (8.4,4.0) {蓝实线：$x^2$\\[8pt]青虚线：$x$\\[8pt]紫点线：$\sqrt{x}$\\[8pt]$0<x<1$：$x^2<x<\sqrt{x}$\\[8pt]$x>1$：$x^2>x>\sqrt{x}$};
\node[text width=12.6cm,align=center] at (7.1,.65) {这里始终固定系数为1；三条曲线在正半轴上只有公共点(1,1)。};
\end{tikzpicture}\end{center}


  \textbf{\textcolor{CStatement}{这个结论解决什么问题？}}

  不必逐个描点，也不必算出具体的幂值，只看指数的范围，就能判断幂函数在正半轴上的走势，并比较两个同底幂的大小；遇到含参数的幂函数，还能先利用系数为 $1$ 的规定求参数。

  \medskip
  \end{statementbox}\clearpage\begin{statementbox}\textbf{\textcolor{CStatement}{条件与定义}}
  \begin{description}[style=nextline,leftmargin=2.8em,labelsep=0.5em,itemsep=0.35em]
    \item[条件1：] 高中所说的幂函数是 $y=x^\alpha$，其中 $\alpha$ 为实数，$x$ 是变量，前面的系数固定为 $1$。
    \item[条件2：] 以下图像走势和大小比较只讨论 $x>0$。此时对任意实数 $\alpha$ 都有 $x^\alpha>0$，图像位于第一象限。
  \end{description}
  \textbf{\textcolor{CStatement}{结论一：先看系数，再看定点}}
  \[
     f(x)=A x^\alpha\ \text{是幂函数}\ \Longrightarrow\ A=1.
  \]
  反过来，若 $A=1$，且指数是实常数、定义域按幂的意义确定，就是高中定义下的幂函数。
  由于 $1^\alpha=1$（包括 $\alpha=0$），所有幂函数都经过 $(1,1)$。而一般的 $y=Ax^\alpha$ 经过 $(1,A)$。若已学习导数，还可得到 $f'(1)=\alpha$，即公共点处的切线斜率等于指数。

  \medskip
  \end{statementbox}\clearpage\begin{statementbox}
\begin{center}\begin{tikzpicture}[font=\small,>=stealth,line join=round]
\draw[gray!40,rounded corners=3pt] (0,0) rectangle (14.2,8.5);\node at (7.1,7.9) {正指数三类：都上升，但陡峭程度不同};
\node at (2.4,6.9) {$\alpha=2$};\begin{scope}[shift={(1.0642857142857143,2.9272727272727272)},x=1.2142857142857144cm,y=0.6545454545454545cm]\draw[->,gray] (-0.3,0)--(2.5,0) node[right]{$x$};\draw[->,gray] (0,-0.5)--(0,5) node[above]{$y$};\draw[blue,thick,solid,domain=0.04:2.1,samples=160] plot(\x,{(\x)^2});\fill[orange!85!black] (1,1) circle(2pt);\end{scope}\node[align=center,text width=4cm] at (2.4,1.7) {$y=x^2$\\上升，越来越陡};\node at (7.050000000000001,6.9) {$\alpha=1$};\begin{scope}[shift={(5.714285714285715,2.9272727272727272)},x=1.2142857142857144cm,y=0.6545454545454545cm]\draw[->,gray] (-0.3,0)--(2.5,0) node[right]{$x$};\draw[->,gray] (0,-0.5)--(0,5) node[above]{$y$};\draw[teal,thick,solid,domain=0.04:2.1,samples=160] plot(\x,{\x});\fill[orange!85!black] (1,1) circle(2pt);\end{scope}\node[align=center,text width=4cm] at (7.050000000000001,1.7) {$y=x$\\斜率始终为1};\node at (11.7,6.9) {$\alpha=\frac12$};\begin{scope}[shift={(10.364285714285714,2.9272727272727272)},x=1.2142857142857144cm,y=0.6545454545454545cm]\draw[->,gray] (-0.3,0)--(2.5,0) node[right]{$x$};\draw[->,gray] (0,-0.5)--(0,5) node[above]{$y$};\draw[violet,thick,solid,domain=0.04:2.1,samples=160] plot(\x,{sqrt(\x)});\fill[orange!85!black] (1,1) circle(2pt);\end{scope}\node[align=center,text width=4cm] at (11.7,1.7) {$y=\sqrt{x}$\\上升，越来越平};
\node[text width=12.6cm,align=center] at (7.1,.65) {三图同尺度，只画正半轴；指数1将两种弯曲走势分开。};
\end{tikzpicture}\end{center}

\textbf{\textcolor{CStatement}{结论二：指数决定正半轴上的五种走势}}

  \begin{center}
  \renewcommand{\arraystretch}{1.35}
  \begin{tabular}{c|c|l}
  指数 & 例子 & 从左向右看图像\\ \hline
  $\alpha>1$ & $x^2$ & 上升，越走越陡\\
  $\alpha=1$ & $x$ & 斜率为 $1$ 的直线\\
  $0<\alpha<1$ & $\sqrt{x}$ & 上升，越走越平\\
  $\alpha=0$ & $1$ & 高度不变\\
  $\alpha<0$ & $1/x$ & 下降，越走越平
  \end{tabular}
  \end{center}
  正指数严格递增，负指数严格递减，零指数为常数。负指数的图像靠近两条坐标轴，但在正半轴上始终不碰到它们。

  \medskip
  \end{statementbox}\clearpage\begin{statementbox}
\begin{center}\begin{tikzpicture}[font=\small,>=stealth,line join=round]
\draw[gray!40,rounded corners=3pt] (0,0) rectangle (14.2,8.8);\node at (7.1,8.200000000000001) {另外两类：零指数水平，负指数下降};
\node at (3.1500000000000004,7.3) {$\alpha=0$，$y=1$};\begin{scope}[shift={(1.2157894736842105,2.88)},x=1.2894736842105265cm,y=0.76cm]\draw[->,gray] (-0.4,0)--(3.4,0) node[right]{$x$};\draw[->,gray] (0,-0.5)--(0,4.5) node[above]{$y$};\draw[gray,thick,solid,domain=0.04:3,samples=180] plot(\x,{1});\fill[orange!85!black] (1,1) circle(2pt);\end{scope}\node[align=center,text width=5.8cm] at (3.1500000000000004,1.7) {每个$x>0$，高度都为1\\导数恒为0};\node at (9.95,7.3) {$\alpha=-1$，$y=1/x$};\begin{scope}[shift={(8.01578947368421,2.88)},x=1.2894736842105265cm,y=0.76cm]\draw[->,gray] (-0.4,0)--(3.4,0) node[right]{$x$};\draw[->,gray] (0,-0.5)--(0,4.5) node[above]{$y$};\draw[orange!85!black,thick,solid,domain=0.25:3,samples=180] plot(\x,{1/\x});\fill[orange!85!black] (1,1) circle(2pt);\end{scope}\node[align=center,text width=5.8cm] at (9.95,1.7) {下降，越来越平\\导数为负，绝对值减小};
\node[text width=12.6cm,align=center] at (7.1,.65) {两图同尺度，只画$x>0$；都经过(1,1)，但不经过原点。};
\end{tikzpicture}\end{center}

\textbf{\textcolor{CStatement}{结论三：底数相同，比较指数}}

  若 $\alpha_1>\alpha_2$，则
  \[
  \begin{cases}
  x^{\alpha_1}>x^{\alpha_2}, & x>1,\\
  x^{\alpha_1}=x^{\alpha_2}=1, & x=1,\\
  x^{\alpha_1}<x^{\alpha_2}, & 0<x<1.
  \end{cases}
  \]
  对 $x>1$ 或 $0<x<1$，上述严格大小关系也能反推指数的大小。

  \medskip
  \textbf{\textcolor{CStatement}{推广结论：两种比较不要混用}}

  指数固定时，比较底数：正指数保持大小方向，负指数颠倒大小方向（底数均为正数）。底数固定时，比较指数：以 $1$ 为分界。两个不同指数的基本幂函数在正半轴上只相交一次，即 $(1,1)$。
\end{statementbox}
